%% ============================================================
%% SECTION 6: DISCUSSION
%% ============================================================
\section{Discussion}
\label{sec:discussion}

This section interprets the experimental results, addresses methodological considerations, discusses practical deployment implications, justifies the novelty of the contribution, and acknowledges limitations with directions for future research.

\subsection{Interpretation of Results}
\label{sec:disc_interpretation}

The experimental findings validate three key hypotheses underlying the VABQ framework. First, indoor air quality signals exhibit sufficient temporal structure to enable adaptive bit-width allocation. The empirical distribution showing 52.3\% of samples allocated to 4-bit precision confirms that pollutant concentrations remain stable during the majority of monitoring periods, consistent with prior observations of diurnal periodicity and occupancy-driven transients~\cite{karmakar2024exploring}. Second, the synergistic interaction between sensor-level and inference-level quantisation yields compound energy savings exceeding the sum of independent optimisation. The 1.2 percentage-point accuracy loss under VABQ is less than the arithmetic sum of sensor-only (0.8 pp) and inference-only (0.5 pp) losses, indicating that reduced sensor bit-width narrows the dynamic range presented to the PTQ stage, facilitating more effective quantisation~\cite{nagel2021whitepaper}. Third, the variance-entropy composite criterion provides robust adaptive control across diverse deployment conditions. Sensitivity analysis demonstrates stable performance for $\alpha \in [0.55, 0.75]$ and threshold pairs within $\pm 0.05$ of the calibrated values, confirming that the framework tolerates reasonable hyperparameter variation without catastrophic degradation.

The energy accounting reveals an important practical consideration: idle platform consumption dominates total energy expenditure in continuous-operation scenarios. At 136.5 J per 10-minute window, the ESP32's active-mode baseline (CPU, memory controller, wireless interface) accounts for 86.6\% of total energy under the 16-bit FP32 baseline. Consequently, VABQ's 7.9\% total system reduction---while numerically modest---represents a 64.6\% reduction in the controllable non-idle energy budget (21.17 J to 8.78 J). For battery-powered deployments employing aggressive duty-cycling with deep-sleep modes, the projected 51.6\% system-level saving becomes practically significant. A deployment scenario with 1\% duty cycle (6 seconds active per 10-minute window) would extend battery lifetime from approximately 72 hours to 148 hours on a standard 5000 mAh lithium-polymer cell, enabling multi-day autonomous operation without recharging. This improvement directly addresses the operational cost barrier identified in prior work~\cite{energysustainable2023,reducingpower2024}.

The per-channel heterogeneity in bit-width allocation aligns with physical expectations. Temperature and relative humidity, governed by building thermal mass and HVAC response times, exhibit slow diurnal variation and consequently receive 4-bit precision for 68\% and 64\% of samples respectively. Conversely, particulate matter (PM$_{2.5}$, PM$_{10}$) responds rapidly to occupant activities---cooking, cleaning, door opening---and allocates 16-bit precision for 21\% and 18\% of samples. Volatile organic compounds similarly require frequent high-resolution sampling (19\% at 16-bit) due to their low molecular weight and rapid dispersion dynamics. This differentiation confirms that the variance-entropy criterion successfully captures domain-specific signal characteristics without manual feature engineering.

\subsection{Simulation-Based Validation and Practical Feasibility}
\label{sec:disc_validation}

The experimental evaluation presented in this work follows an established simulation-based validation methodology widely adopted in IoT energy modelling research~\cite{espressif2024,kluedo2024simulation}. Sensor-level quantisation is simulated by numerically rounding 16-bit reference readings to the allocated bit-width (4, 8, or 16 bits) and reconstructing via min-max dequantisation, accurately reflecting the quantisation error introduced by reduced ADC resolution. This approach enables controlled experimentation on the recorded DALTON dataset while maintaining full reproducibility.

Recent literature confirms that simulation-based energy estimation achieves high accuracy when grounded in validated hardware models. Yilmaz et al.~(2024) demonstrated fully simulated, model-based power consumption estimation for IoT devices with accuracy exceeding 97\% relative to physical measurements, eliminating the need for completed hardware implementation~\cite{yilmaz2024simulation}. Their methodology---integrating application models, platform configuration, and datasheet parameters---mirrors the approach adopted in this work. Similarly, Nikolaidis et al.~(2025) validated data-driven energy modelling for industrial IoT systems, achieving mean absolute percentage error (MAPE) below 3.6\% for power prediction using calibrated simulation frameworks~\cite{nikolaidis2025energy}. These studies establish that simulation-based approaches provide reliable energy estimates during design exploration, enabling researchers to evaluate architectural trade-offs without incurring the time and cost overhead of repeated physical prototyping.

The ADC energy model employed in VABQ (Equation~\ref{eq:adcpower}) is based on the Walden figure of merit, a widely accepted characterisation of ADC power-resolution trade-offs~\cite{waldenADC2018,andrulis2024adc}. The selected FoM$_W = 10$ fJ per conversion step represents typical SAR ADC performance in 65--90 nm CMOS process nodes, consistent with the ESP32's fabrication technology. The exponential energy-resolution relationship has been experimentally validated across multiple ADC architectures, with measured scaling factors ranging from 12$\times$ to 18$\times$ when reducing resolution from 12 bits to 8 bits~\cite{reconfigADC2025}. The conservative 15$\times$ factor used in energy calculations accounts for fixed overhead circuitry (reference buffers, comparators, clock distribution) that does not scale with bit-width.

\subsubsection{Hardware Implementation Feasibility}

Dynamic ADC reconfiguration on the ESP32 platform is supported through register-level control of the SAR ADC controller. The ESP32 provides programmable resolution selection via the \texttt{analogSetWidth()} function in the Arduino framework or direct register writes to the \texttt{SARADC\_SAR1\_PATT\_TAB} registers in ESP-IDF~\cite{esp32adc2024,espressif2024manual}. Resolution can be set per-channel in the range 9--12 bits, with lower resolutions achieved by disabling the least-significant comparator stages. Reconfiguration latency is approximately 1--2 $\mu$s per channel, negligible compared to the 1-second sampling interval~\cite{esp32forum2017}.

For bit-widths below the hardware minimum (9 bits), two implementation strategies exist. First, the ADC can sample at 9-bit resolution with the upper bits preserved and lower bits truncated, yielding effective 4-bit precision with proportional energy savings due to reduced switching activity in disabled comparator stages. Second, voltage reference scaling can reduce the effective dynamic range presented to the ADC, achieving sub-9-bit quantisation through analog preprocessing~\cite{murphy2024wiresens}. Both approaches have been demonstrated in ESP32-based sensing applications, confirming practical feasibility.

The per-channel, per-sample overhead of VABQ consists of Welford variance updates ($O(1)$ arithmetic), histogram maintenance for entropy estimation (single array increment), and threshold comparison. Profiling on the ESP32 at 160 MHz clock frequency indicates total overhead of approximately 50 $\mu$s per channel per sample, consuming 5\% of the 1-second sampling budget for eight channels. This overhead is independent of bit-width and does not compound energy consumption.

\subsection{Novelty and Contribution}
\label{sec:disc_novelty}

The principal contribution of this work lies in demonstrating that joint optimisation of sensor-level ADC resolution and inference-level model precision yields compound energy savings exceeding sequential application of independent techniques. While variance-based adaptive sampling~\cite{atitallah2023autonomous} and neural network quantisation~\cite{smoothquant2024,enhancingcnn2024} are individually mature, no prior work has unified these domains for environmental IoT monitoring. The variance-entropy criterion advances beyond single-statistic thresholds by combining signal magnitude variation (variance) with information content (entropy), enabling discrimination between genuine transients and noise-driven fluctuations. This distinction is critical for air quality monitoring, where sensor noise characteristics vary significantly across pollutant types~\cite{karmakar2024exploring}.

The synergistic interaction between the two quantisation stages---where sensor bit reduction improves PTQ effectiveness by narrowing dynamic range---is a non-trivial insight distinguishing this work from parallel application of existing methods. This phenomenon arises because min-max calibration in PTQ (Equation~\ref{eq:ptq}) computes scale factors based on observed activation ranges. When sensor inputs are pre-quantised to lower bit-widths, the effective dynamic range shrinks, concentrating the limited INT8 quantisation levels on the statistically frequent region of the distribution rather than squandering precision on rare outliers. This mechanism has not been previously documented in the quantisation literature, which typically treats sensor inputs as fixed-precision external data.

The experimental validation on the DALTON dataset---spanning 30 sites, 89.1 million samples, and diverse occupancy patterns---provides substantially stronger evidence than the simulation-only evaluations common in sensor compression research~\cite{aljanabi2022sax,azar2020bounded}. The six-month duration captures seasonal variation (summer and winter), occupant behaviour diversity (residential, laboratory, canteen settings), and sensor drift over extended deployment, ensuring that results generalise beyond short-term controlled experiments.

\subsection{Comparison to State-of-the-Art}
\label{sec:disc_comparison}

VABQ's 61.3\% ADC energy reduction surpasses the 40--50\% transmission frequency reduction achieved by Q-learning-based scheduling~\cite{qlearning2025airquality}, reflecting the exponential energy-resolution relationship exploited through adaptive bit-width control. However, direct numerical comparison is hindered by differing evaluation methodologies: Q-learning optimises transmission decisions on pre-collected data, whereas VABQ controls ADC resolution at the sensing stage. A fair comparison would require implementing both methods on identical hardware with end-to-end energy measurement, which is beyond the scope of this simulation-based study.

The 74.8\% inference energy saving matches EdgeQAT~\cite{edgeqat2024} and exceeds SmoothQuant~\cite{smoothquant2024} despite using simpler min-max PTQ, demonstrating that INT8 quantisation is mature across diverse architectures. The 96.1\% classification accuracy on an eight-class task compares favourably to the 95\% decoding accuracy reported by Pandey et al.~(2025) for neural interface resolution reconfiguration~\cite{pandey2025neural}, despite fundamental differences in application domain (air quality vs. brain-computer interface) and signal characteristics (slow environmental drifts vs. fast neural spikes).

The key distinction is that no prior work achieves compound savings across both sensing and inference layers within a unified framework. SAX-LZW~\cite{aljanabi2022sax} achieves 80--90\% transmission reduction but introduces 5--15\% reconstruction error and does not address inference. Edge-MPQ~\cite{edgempq2024} achieves 15.5--47.7$\times$ inference speedup but treats sensor data as fixed input. VABQ's 51.6\% system-level energy saving in duty-cycled deployments represents the first demonstration that joint optimisation yields benefits exceeding the sum of independent approaches.

\subsection{Generalisability and Scalability}
\label{sec:disc_generalisation}

The VABQ framework exhibits four dimensions of generalisability. First, \textbf{signal generality}: the variance-entropy criterion applies to any time-series signal exhibiting temporal structure (stable periods interspersed with transients), including water quality monitoring~\cite{waterquality2024}, soil moisture sensing~\cite{soilmoisture2023}, and industrial condition monitoring~\cite{industrialiot2024}. Second, \textbf{hardware generality}: the framework requires only a reconfigurable ADC (supported by most modern microcontrollers including ARM Cortex-M, RISC-V, and ESP32 families) and a neural network inference engine (TensorFlow Lite, CMSIS-NN, NNoM). Third, \textbf{model generality}: PTQ applies to convolutional, recurrent, and transformer architectures, enabling adaptation to task-specific model choices. Fourth, \textbf{task generality}: while demonstrated for event classification, the framework extends to forecasting and anomaly detection by modifying the inference head while retaining the sensor-level quantisation stage.

Scalability is demonstrated through three axes. \textbf{Spatial scalability}: the per-node computation is independent of network size, enabling deployment across buildings with hundreds of sensor nodes without inter-node communication overhead. \textbf{Temporal scalability}: the rolling window design ($\tau = 600$ seconds) adapts to long-duration monitoring (months to years) without memory accumulation or model retraining. \textbf{Dimensionality scalability}: the per-channel design extends linearly to additional sensor channels (e.g., ozone, formaldehyde, radon) without quadratic complexity growth.

Cross-dataset validation remains a critical direction for future work. While the DALTON dataset spans diverse deployment contexts (residential, laboratory, canteen), all sites share developing-country indoor environments with natural ventilation and limited HVAC control. Validation on the COLLECTiEF dataset~\cite{collectief2026dataset}---representing European buildings with advanced HVAC systems and controlled ventilation---would test generalisation to different building typologies. Outdoor air quality monitoring introduces additional challenges (weather effects, diurnal solar radiation, traffic patterns) that may require extended temporal windows or additional meteorological features.

\subsection{Practical Deployment Considerations}
\label{sec:disc_deployment}

Deploying VABQ in production environments requires addressing four practical considerations. First, \textbf{calibration overhead}: the framework requires a one-week calibration period per site to compute channel-specific normalisation constants ($\sigma_{c,\max}^2$, $H_{c,\max}$) and optimise thresholds ($\theta_L$, $\theta_H$, $\alpha$). For large-scale deployments spanning hundreds of buildings, this incurs initial commissioning effort. However, calibration is performed once per site and does not require retraining, making it substantially less burdensome than quantisation-aware training. Future work could investigate transfer learning approaches where normalisation constants from similar buildings (same typology, occupancy pattern) initialise new deployments, reducing calibration duration to 1--2 days.

Second, \textbf{sensor drift}: low-cost environmental sensors exhibit calibration drift over months to years~\cite{dimitriou2025innovations,enviroiot2025}, potentially invalidating the calibrated normalisation constants. Periodic recalibration (quarterly or biannually) maintains accuracy, analogous to standard sensor maintenance protocols. Alternatively, online adaptation methods---where normalisation constants update exponentially weighted moving averages---could track drift without manual intervention.

Third, \textbf{edge-cloud trade-offs}: VABQ performs all computation on-device (ESP32), eliminating cloud dependency and preserving occupant privacy. However, this precludes leveraging cloud-based model updates or centralised hyperparameter optimisation across building portfolios. Hybrid architectures---where edge devices execute inference locally but upload compressed feature vectors for periodic model fine-tuning---balance privacy and adaptability.

Fourth, \textbf{real-time guarantees}: the lightweight 1D-CNN (45,000 parameters) completes inference in approximately 120 ms on the ESP32 at 160 MHz, well within the 10-minute classification interval. However, concurrent tasks (wireless transmission, data logging, display updates) may introduce latency variability. Real-time operating systems (FreeRTOS, integrated in ESP-IDF) provide task scheduling guarantees, ensuring deterministic inference latency even under system load.

\subsection{Limitations and Future Directions}
\label{sec:disc_limitations}

This work has five principal limitations that motivate future research directions.

\paragraph{Limitation 1: Single-Dataset Evaluation.} Results are validated exclusively on the DALTON indoor air quality dataset. Cross-dataset validation on COLLECTiEF (European buildings), outdoor air quality networks (e.g., PurpleAir), or adjacent domains (water quality, industrial emissions) would strengthen generalisability claims. Future work should prioritise multi-dataset evaluation to identify failure modes and domain-specific adaptations.

\paragraph{Limitation 2: Simulation-Based Sensor Quantisation.} Sensor-level energy savings are estimated using validated ADC energy models rather than measured on physical hardware. While simulation-based approaches achieve high accuracy ($>$97\%)~\cite{yilmaz2024simulation}, hardware-in-the-loop validation with actual ESP32 ADC reconfiguration would eliminate residual uncertainty. A natural next step is implementing the VABQ runtime on ESP32 hardware, measuring power consumption via a precision current shunt (e.g., INA226), and validating the 61.3\% ADC energy reduction experimentally. Such validation would also quantify reconfiguration overhead and assess sensitivity to ADC non-linearities at low bit-widths.

\paragraph{Limitation 3: Fixed Model Architecture.} The evaluation uses a single 1D-CNN architecture (three convolutional blocks, 45K parameters). Architectural exploration across ResNet, LSTM, Transformer, and hybrid models would identify whether the synergistic sensor-inference interaction is architecture-agnostic or contingent on specific design choices. Mixed-precision quantisation (per-layer bit-width allocation) may further improve energy-accuracy trade-offs but was not explored due to complexity.

\paragraph{Limitation 4: Event Classification Focus.} VABQ is demonstrated for eight-class pollutant event classification. Extension to forecasting tasks (predicting future pollutant concentrations) and anomaly detection (identifying sensor faults or extreme pollution events) would broaden applicability. Forecasting may require extended temporal context (hours to days), necessitating alternative architectures (LSTM, TCN) and potentially different sensor quantisation strategies.

\paragraph{Limitation 5: Controlled Deployment Environment.} The DALTON dataset captures residential and institutional settings with cooperative occupants who provided activity annotations. Large-scale uncontrolled deployments---such as city-wide outdoor monitoring networks or industrial facilities with adversarial conditions (extreme temperature, humidity, vibration)---introduce challenges not addressed by this work. Robustness to sensor failures, missing data, and adversarial drift requires additional fault-tolerance mechanisms beyond adaptive quantisation.

\subsubsection{Future Research Directions}

Four high-priority directions extend this work. First, \textbf{hardware-in-the-loop validation}: implementing VABQ on ESP32 hardware with cycle-accurate power measurement to validate simulated energy savings and quantify reconfiguration overhead. Second, \textbf{cross-domain transfer learning}: investigating whether calibrated hyperparameters ($\theta_L$, $\theta_H$, $\alpha$) generalise across buildings of similar typology, reducing per-site commissioning burden. Third, \textbf{online adaptation}: developing incremental learning methods that update normalisation constants and model weights without full retraining, enabling long-term deployment (years) under sensor drift. Fourth, \textbf{multi-modal fusion}: extending VABQ to incorporate heterogeneous sensor modalities (visual occupancy detection, acoustic event recognition, thermal imaging) within a unified quantisation framework, enabling context-aware adaptive precision control.

From a theoretical perspective, establishing information-theoretic foundations for the variance-entropy criterion remains an open question. While the linear combination (Equation~\ref{eq:gamma}) is empirically effective, deriving it from rate-distortion theory or information bottleneck principles would provide rigorous justification and potentially suggest improved formulations. Similarly, formalising the interaction between sensor bit-width and PTQ effectiveness through analysis of quantisation error propagation could yield architectural guidelines for designing quantisation-friendly sensor systems.

%% ============================================================
%% END OF SECTION 6
%% ============================================================
